% ============================================================
%  Pipeline Híbrido LLM + ESBMC — Qualificação de Mestrado
%  Esqueleto — Introdução detalhada, demais seções em rascunho
%  Última atualização: 2026-06-05
% ============================================================
\documentclass[conference]{IEEEtran}
\IEEEoverridecommandlockouts

\usepackage[utf8]{inputenc}
\usepackage[T1]{fontenc}
\usepackage[brazil]{babel}
\usepackage{cite}
\usepackage{amsmath,amssymb,amsfonts}
\usepackage{algorithmic}
\usepackage{algorithm}
\usepackage{graphicx}
\usepackage{textcomp}
\usepackage{xcolor}
\usepackage{booktabs}
\usepackage{hyperref}
\usepackage{listings}

\def\BibTeX{{\rm B\kern-.05em{\sc i\kern-.025em b}\kern-.08em
    T\kern-.1667em\lower.7ex\hbox{E}\kern-.125emX}}

% ── Macros de conveniência ───────────────────────────────────
\newcommand{\esbmc}{\textsc{Esbmc}}
\newcommand{\flowA}{Flow~A}
\newcommand{\flowB}{Flow~B}
\newcommand{\flowC}{Flow~C}
\newcommand{\todo}[1]{\textcolor{red}{\textbf{[TODO: #1]}}}

% ============================================================
\begin{document}

\title{Verificação Formal Guiada por LLM para Detecção de Bugs de
       Runtime em Código Python: Um Pipeline Híbrido com \esbmc{}
       \thanks{%
         \todo{Identificar agência de fomento.}}}

\author{%
  \IEEEauthorblockN{\todo{Nome Completo}}
  \IEEEauthorblockA{%
    \textit{\todo{Departamento}}\\
    \textit{\todo{Universidade}}\\
    \todo{Cidade, País}\\
    \todo{e-mail}}
  \and
  \IEEEauthorblockN{\todo{Orientador}}
  \IEEEauthorblockA{%
    \textit{\todo{Departamento}}\\
    \textit{\todo{Universidade}}\\
    \todo{Cidade, País}\\
    \todo{e-mail}}}

\maketitle

% ── Abstract ────────────────────────────────────────────────
\begin{abstract}
A verificação formal de código Python apresenta desafios práticos
significativos: verificadores baseados em \textit{Bounded Model Checking}
(BMC) requerem um ponto de entrada explícito para analisar funções
isoladas, enquanto Modelos de Linguagem de Grande Escala (LLMs) identificam
bugs com alta revocação mas produzem falsos positivos não verificados.
Este trabalho propõe um pipeline híbrido de três fluxos que combina análise
semântica por LLM com verificação formal pelo \esbmc{} para detectar bugs
de runtime em código Python.
No \flowB{} (principal), a LLM propõe hipóteses de bug identificando a
função suspeita e a categoria do defeito; o \esbmc{} as verifica
formalmente usando o flag \texttt{--function}, que torna os parâmetros da
função simbólicos sem instrumentação manual.
O \flowA{} (ESBMC-\textit{only}) e o \flowC{} (LLM-\textit{only}) servem
como \textit{baselines} para isolar a contribuição de cada componente.
Avaliamos os três fluxos em um dataset de 70 arquivos Python anotados
manualmente, cobrindo três categorias de bugs formais
(\textit{assertion\_violation}, \textit{division\_by\_zero},
\textit{out\_of\_bounds}) e três categorias de \textit{code smells}.
\todo{Preencher com resultados obtidos após execução do benchmark.}
\end{abstract}

\begin{IEEEkeywords}
verificação formal, bounded model checking, ESBMC, large language model,
detecção de bugs, Python, pipeline híbrido, análise estática
\end{IEEEkeywords}

% ============================================================
\section{Introdução}
\label{sec:introducao}

% ── 1.1 Contexto ────────────────────────────────────────────
\subsection{Contexto}
\label{sec:contexto}

Python consolidou-se como uma das linguagens de programação mais utilizadas
globalmente, presente em domínios que vão de ciência de dados e aprendizado
de máquina a desenvolvimento web e automação de infraestrutura
\todo{citar índice TIOBE/GitHub Octoverse}.
Apesar de sua produtividade, a natureza dinamicamente tipada da linguagem
torna a detecção estática de erros de runtime — como
\texttt{IndexError}, \texttt{ZeroDivisionError} e \texttt{AssertionError}
— uma tarefa particularmente desafiadora
\todo{citar trabalho sobre bugs comuns em Python}.

Verificação formal baseada em \textit{Bounded Model Checking} (BMC)
oferece garantias matemáticas: dado um bound $k$, se nenhuma violação
de propriedade é encontrada, o programa está livre de erros até
profundidade $k$~\todo{citar Clarke et al., Biere et al.}.
O \esbmc{}~\todo{citar ESBMC paper principal} é um verificador BMC maduro
com suporte a código Python via \texttt{--python python3}, capaz de gerar
contraexemplos concretos para propriedades como \textit{array bounds},
divisão por zero e violação de asserções.

Simultaneamente, Modelos de Linguagem de Grande Escala (LLMs) demonstraram
capacidade de identificar padrões de bugs em código-fonte de forma
semântica, sem instrumentação~\todo{citar CodeLlama, GPT-4 code review,
DeepCode}.
Entretanto, LLMs produzem falsos positivos (\textit{alucinações}): reportam
bugs em expressões que não existem no código executável, ou classificam
incorretamente guardas existentes~\todo{citar trabalho sobre alucinações
em LLM code review}.

% ── 1.2 Problema ────────────────────────────────────────────
\subsection{Problema}
\label{sec:problema}

Duas limitações complementares impedem o uso isolado de cada abordagem:

\textbf{Limitação do BMC:} O \esbmc{} e verificadores similares exigem
um ponto de entrada explícito — tipicamente uma função \texttt{main} ou
um \textit{harness} instrumentado manualmente — para verificar funções
de biblioteca isoladas.
Em bases de código Python reais, funções utilitárias não possuem
\texttt{main}; construir harnesses por função não escala para análise
de repositórios.
O flag \texttt{--function} do \esbmc{} resolve parcialmente esse problema:
ao especificar \texttt{--function~<nome>}, os parâmetros da função tornam-se
variáveis simbólicas automaticamente, permitindo verificação sem harness.
Contudo, o flag requer que o analista saiba \emph{qual} função verificar
e \emph{qual} propriedade testar — informações que o BMC por si só
não pode inferir.

\textbf{Limitação da LLM:} LLMs identificam a função suspeita e a
categoria do potencial bug com alta revocação, mas sem garantia formal.
Falsos positivos da LLM — reportar bugs em expressões protegidas por
guardas corretas, ou citar expressões inexistentes — consomem tempo de
verificação e reduzem a confiança do desenvolvedor.

A Figura~\ref{fig:gap} ilustra a lacuna que ambas as abordagens deixam
quando operadas de forma independente.

\begin{figure}[ht]
  \centering
  % \includegraphics[width=0.9\columnwidth]{figs/gap_diagram.pdf}
  \todo{Inserir diagrama: LLM (alta revocação, sem prova) vs ESBMC
        (prova formal, sem alvo). Espaço intermediário = pipeline híbrido.}
  \caption{Lacuna entre análise semântica por LLM e verificação formal
           por BMC para funções Python isoladas.}
  \label{fig:gap}
\end{figure}

% ── 1.3 Motivação / Justificativa ───────────────────────────
\subsection{Motivação e Justificativa}
\label{sec:motivacao}

A combinação dessas duas abordagens segue o princípio
\emph{``LLM Proposes, Formal Tools Dispose''}:
a LLM atua como oráculo semântico para identificar a função-alvo e
a hipótese de bug; o \esbmc{} atua como árbitro formal para confirmar
ou refutar a hipótese com um contraexemplo concreto ou um certificado
de segurança dentro do bound.

Esta divisão de responsabilidades oferece três vantagens:
\begin{enumerate}
  \item \textbf{Escalabilidade da verificação formal:} O \esbmc{} só é
    invocado para funções e propriedades pré-filtradas pela LLM, reduzindo
    o espaço de busca.
  \item \textbf{Eliminação de alucinações:} Findings da LLM que não
    encontram confirmação formal são explicitamente classificados como
    \texttt{not\_confirmed\_within\_bound} ou \texttt{llm\_false\_positive},
    separando suspeitas de provas.
  \item \textbf{Auditabilidade:} O contraexemplo gerado pelo \esbmc{}
    constitui evidência reproducível da presença do bug, independente
    do modelo de linguagem utilizado.
\end{enumerate}

Adicionalmente, a comparação sistemática dos três fluxos no mesmo dataset
permite isolar a contribuição de cada componente:
o \flowA{} (ESBMC-\textit{only}) estabelece o teto alcançável sem
semântica; o \flowC{} (LLM-\textit{only}) estabelece a linha base sem
prova formal; o \flowB{} (híbrido) é a proposta principal.

% ── 1.4 Objetivos ───────────────────────────────────────────
\subsection{Objetivos}
\label{sec:objetivos}

\textbf{Objetivo geral:} Propor, implementar e avaliar um pipeline
híbrido que combina análise semântica por LLM com verificação formal
pelo \esbmc{} para detectar bugs de runtime em funções Python, comparando
sistematicamente os três fluxos de execução em um dataset anotado.

\textbf{Objetivos específicos:}
\begin{enumerate}
  \item[\textbf{OE1.}] Construir um dataset de 70 arquivos Python com
    anotação manual de ground truth por função, cobrindo três categorias
    de bugs formalmente verificáveis e três categorias de \textit{code
    smells} (\S\ref{sec:dataset}).
  \item[\textbf{OE2.}] Implementar o pipeline de três fluxos
    (Flow~A / Flow~B / Flow~C) e o framework de avaliação com métricas
    de precisão, revocação e F1 por fluxo e por categoria
    (\S\ref{sec:metodo}).
  \item[\textbf{OE3.}] Responder às seguintes questões de pesquisa:
    \begin{description}
      \item[\textbf{RQ1:}] O pipeline híbrido (\flowB{}) alcança F1
        superior ao \esbmc{}-\textit{only} (\flowA{}) na detecção de
        bugs formais em Python?
      \item[\textbf{RQ2:}] O pipeline híbrido (\flowB{}) alcança F1
        superior ao LLM-\textit{only} (\flowC{}) em precisão formal?
      \item[\textbf{RQ3:}] Qual categoria de bug (\textit{assertion
        violation}, \textit{division by zero}, \textit{out of bounds})
        mais beneficia da combinação LLM + \esbmc{}?
    \end{description}
  \item[\textbf{OE4.}] Quantificar a taxa de alucinação da LLM antes
    e depois da confirmação formal (\S\ref{sec:avaliacao}).
\end{enumerate}

O restante deste documento está organizado da seguinte forma:
a Seção~\ref{sec:correlatos} discute os trabalhos relacionados;
a Seção~\ref{sec:referencial} apresenta o referencial teórico;
a Seção~\ref{sec:metodo} descreve o método proposto;
a Seção~\ref{sec:resultados} apresenta os resultados parciais;
a Seção~\ref{sec:metodologia} detalha o protocolo experimental e
o cronograma;
a Seção~\ref{sec:conclusao} conclui o trabalho.

% ============================================================
\section{Trabalhos Correlatos}
\label{sec:correlatos}

% ── O que buscar na literatura ──────────────────────────────
% Grupos a cobrir:
%   (a) LLM para análise e detecção de bugs em código:
%       CodeLlama, Copilot, CodeT5, DeepCode, VulBERTa, etc.
%   (b) LLM + verificação formal (híbridos):
%       LLM para síntese de invariantes (Houdini+LLM),
%       LLM para geração de harnesses (FuzzGPT, etc.)
%       "LLM Proposes, Formal Tools Dispose" (citar se existir)
%   (c) ESBMC e BMC para Python:
%       Papers do grupo ESBMC sobre suporte Python,
%       PyExZ3, Pyre, mypy (análise estática Python)
%   (d) Datasets de bugs Python:
%       BugsInPy, PyBugs, etc.
%
% Diferencial deste trabalho:
%   - Usa --function (sem harness manual) para Python
%   - Compara 3 fluxos sistematicamente no mesmo dataset
%   - Mede alucinação LLM via validação AST

\todo{Preencher com revisão bibliográfica sistemática.
      Grupos: LLM code review, LLM+formal, ESBMC Python, datasets.
      Destacar como o pipeline híbrido com --function se diferencia.}

% ============================================================
\section{Referencial Teórico}
\label{sec:referencial}

\subsection{Bounded Model Checking (BMC)}
\todo{Definir BMC formalmente: fórmula SAT, bound k, contraexemplo.
      Citar Clarke, Biere, CBMC.}

\subsection{ESBMC e Verificação de Código Python}
\todo{Descrever flag --function: parâmetros simbólicos automáticos.
      Limitações: sem len(), subset de Python suportado.
      Citar papers ESBMC.}

\subsection{Modelos de Linguagem de Grande Escala para Código}
\todo{GPT-4 / Responses API. Structured outputs (JSON schema).
      Temperature=0 para reproducibilidade.
      Conceito de alucinação em code review.}

\subsection{Análise Estática via AST em Python}
\todo{Papel da validação AST para filtrar falsos positivos da LLM.
      Módulo ast do Python.}

% ============================================================
\section{Método Proposto}
\label{sec:metodo}

% Este seção é 100% derivável do repositório — sem dependência de bib.

\subsection{Visão Geral do Pipeline}

O pipeline é composto por cinco estágios sequenciais, descritos a seguir
e ilustrados na Figura~\ref{fig:pipeline}:

\begin{enumerate}
  \item \textbf{Preprocessamento (AST):} extração de \texttt{CodeUnit}
    por função — parâmetros, operações perigosas (divisões, subscripts,
    asserts), guardas e métricas de complexidade.
  \item \textbf{Análise LLM:} para cada \texttt{CodeUnit}, o modelo LLM
    recebe um prompt estruturado com 8 passos de raciocínio obrigatórios
    e retorna \texttt{Finding}s em JSON validado por schema.
  \item \textbf{Validação AST:} findings com \texttt{verifiable=true} têm
    a expressão reportada verificada no AST executável; expressões
    ausentes são classificadas como \texttt{llm\_false\_positive}.
  \item \textbf{Verificação ESBMC (\flowB{}):} findings verificáveis são
    submetidos ao \esbmc{} com \texttt{--function <nome>} e flags
    específicos da categoria (Tabela~\ref{tab:flags}).
  \item \textbf{Consolidação:} cada finding recebe uma classificação
    final determinística (Tabela~\ref{tab:classificacoes}).
\end{enumerate}

\begin{figure}[ht]
  \centering
  % \includegraphics[width=0.95\columnwidth]{figs/pipeline.pdf}
  \todo{Inserir diagrama de fluxo do pipeline: AST → LLM → Validação
        AST → ESBMC --function → Classificação final.}
  \caption{Arquitetura do pipeline híbrido LLM + \esbmc{}.}
  \label{fig:pipeline}
\end{figure}

\subsection{Flow A — ESBMC-\textit{Only}}

No \flowA{}, o \esbmc{} é invocado com \texttt{--function~<fn>} para
cada função extraída pelo preprocessador, sem envolvimento da LLM.
Os parâmetros tornam-se simbólicos automaticamente; nenhum harness é
necessário.
O resultado serve como \textit{baseline} formal — teto de desempenho
alcançável por BMC sem orientação semântica.

\subsection{Flow B — Pipeline Híbrido (Principal)}

No \flowB{}, a LLM recebe o código da função junto com as operações
detectadas pelo preprocessador e retorna hipóteses de bug em JSON.
Para cada hipótese com \texttt{verifiable=true}, o \esbmc{} é invocado
com \texttt{--function} e os flags da categoria (Tabela~\ref{tab:flags}).

\begin{table}[htbp]
  \caption{Flags \esbmc{} por categoria de bug no \flowB{}}
  \label{tab:flags}
  \begin{center}
    \begin{tabular}{ll}
      \toprule
      \textbf{Categoria} & \textbf{Flags adicionais} \\
      \midrule
      \texttt{division\_by\_zero}    & \texttt{--no-bounds-check} \\
      \texttt{out\_of\_bounds}       & \texttt{--no-div-by-zero-check --assign-param-nondet} \\
      \texttt{assertion\_violation}  & (nenhum) \\
      \bottomrule
    \end{tabular}
  \end{center}
\end{table}

\subsection{Flow C — LLM-\textit{Only} (Baseline Neural)}

O \flowC{} é idêntico ao \flowB{} até a etapa de análise LLM, mas omite
a verificação pelo \esbmc{}.
Findings com \texttt{verifiable=true} são classificados como
\texttt{llm\_only\_suspected}, sem confirmação formal.
Serve como \textit{baseline} neural — mede a capacidade da LLM antes
da triagem formal.

\subsection{Classificações Finais}

Cada finding recebe uma classificação determinística descrita na
Tabela~\ref{tab:classificacoes}.

\begin{table}[htbp]
  \caption{Classificações finais do pipeline}
  \label{tab:classificacoes}
  \begin{center}
    \begin{tabular}{ll}
      \toprule
      \textbf{Classificação} & \textbf{Significado} \\
      \midrule
      \texttt{llm\_confirmed\_by\_esbmc}   & LLM + \esbmc{} confirmaram (TP híbrido) \\
      \texttt{esbmc\_native\_bug}           & \flowA{} detectou, LLM não havia apontado \\
      \texttt{llm\_missed\_esbmc\_bug}      & FN da LLM detectado pelo \flowA{} \\
      \texttt{not\_confirmed\_within\_bound}& \esbmc{} verificou sem violação no bound \\
      \texttt{esbmc\_inconclusive}          & Timeout, erro ou resultado ambíguo \\
      \texttt{llm\_false\_positive}         & Expressão ausente no AST (alucinação) \\
      \texttt{llm\_only\_suspected}         & \flowC{}: suspeita sem verificação \\
      \texttt{heuristic\_smell\_only}       & Code smell, sem verificação formal \\
      \texttt{skipped\_not\_verifiable}     & Fora do escopo do \flowB{} atual \\
      \texttt{out\_of\_scope\_finding}      & Categoria não suportada na V1 \\
      \bottomrule
    \end{tabular}
  \end{center}
\end{table}

\subsection{Dataset}
\label{sec:dataset}

O dataset é composto por \todo{N} arquivos Python construídos manualmente,
distribuídos conforme a Tabela~\ref{tab:dataset}.
Todos os arquivos foram validados para não conterem chamadas a
\texttt{len()} (limitação do suporte \esbmc{} a Python).
O ground truth é anotado em JSON por função, especificando categoria
esperada, verifiability, e expressão envolvida.

\begin{table}[htbp]
  \caption{Distribuição do dataset por categoria}
  \label{tab:dataset}
  \begin{center}
    \begin{tabular}{llr}
      \toprule
      \textbf{Tipo} & \textbf{Categoria} & \textbf{Arquivos} \\
      \midrule
      Bug formal & \texttt{assertion\_violation}  & 15 \\
      Bug formal & \texttt{division\_by\_zero}    & 15 \\
      Bug formal & \texttt{out\_of\_bounds}       & 15 \\
      Controle   & \texttt{clean}                 & 10 \\
      Smell      & \texttt{complex\_conditional}  & 5  \\
      Smell      & \texttt{long\_method}          & 5  \\
      Smell      & \texttt{many\_parameters}      & 5  \\
      \midrule
                 & \textbf{Total}                 & \textbf{70} \\
      \bottomrule
    \end{tabular}
  \end{center}
\end{table}

\subsection{Avaliação}
\label{sec:avaliacao}

Para cada fluxo, calculamos precisão ($P$), revocação ($R$) e $F_1$:

\begin{equation}
  P = \frac{TP}{TP + FP}, \quad
  R = \frac{TP}{TP + FN}, \quad
  F_1 = \frac{2PR}{P + R}
  \label{eq:prf}
\end{equation}

Bugs formais e \textit{code smells} são avaliados separadamente.
Arquivos \texttt{clean} são controles negativos: qualquer finding
reportado conta como FP.
A taxa de alucinação é calculada como:

\begin{equation}
  \text{Hal} = \frac{\texttt{llm\_false\_positive}}
               {\texttt{bug\_tp} + \texttt{bug\_fp} +
                \texttt{llm\_false\_positive}}
  \label{eq:hal}
\end{equation}

% ============================================================
\section{Resultados Parciais}
\label{sec:resultados}

\todo{Esta seção será preenchida após execução do benchmark completo.
      Rodar: python scripts/smoke\_benchmark.py --model gpt-4.1-nano --n 15 --save-json artifacts/results/flow\_ab\_full.json
      e o equivalente para Flow C.}

% Estrutura planejada das tabelas de resultados:
%
% Tabela A — P/R/F1 por fluxo (Flow A / Flow B / Flow C)
% Tabela B — F1 por categoria de bug (av / dz / oob) × fluxo
% Tabela C — Resultados da LLM: confirmed / not confirmed / FP / inconclusive
% Tabela D — Taxa de alucinação por modelo

% ============================================================
\section{Metodologia e Cronograma}
\label{sec:metodologia}

\subsection{Protocolo Experimental}

Os três fluxos são avaliados sobre o mesmo dataset (70 arquivos),
com os parâmetros fixos: bound $k = 5$, timeout $= 30$s,
temperatura LLM $= 0$.
Cada arquivo é processado independentemente; não há contaminação
cruzada entre fluxos.
O modelo LLM principal é \todo{gpt-4.1-nano} (ou outro a especificar).

Para comparação entre modelos, o script aceita \texttt{--model} como
parâmetro, permitindo rodar o mesmo protocolo com modelos diferentes
e gerar JSONs independentes consolidados no dashboard HTML.

\subsection{Cronograma}

\todo{Inserir tabela de cronograma de atividades restantes:
      coleta de resultados, revisão bibliográfica, redação, submissão.}

% ============================================================
\section{Conclusão}
\label{sec:conclusao}

\todo{Preencher após obtenção dos resultados. Estrutura sugerida:
      (i) retomar hipótese, (ii) responder RQ1/RQ2/RQ3,
      (iii) limitações, (iv) trabalhos futuros.}

A abordagem híbrida proposta demonstra que o problema de ponto de entrada
do BMC em código Python pode ser resolvido sem harness manual, usando
o flag \texttt{--function} do \esbmc{} direcionado pela LLM.
A separação em três fluxos independentes (\flowA{}, \flowB{}, \flowC{})
permite atribuir ganhos de desempenho a cada componente de forma
controlada.
O framework de 10 classificações finais e as métricas de alucinação
oferecem uma visão granular do comportamento do pipeline que vai além
das métricas agregadas de F1.

% ============================================================
\section*{Agradecimentos}

\todo{Preencher com agências de fomento e colaboradores.}

% ============================================================
\begin{thebibliography}{00}

% ── BMC / ESBMC ─────────────────────────────────────────────
\bibitem{esbmc}
  \todo{Referência principal do ESBMC.}
\bibitem{bmc_clarke}
  \todo{Clarke et al. — Bounded Model Checking.}
\bibitem{biere_bmc}
  \todo{Biere et al. — BMC survey.}

% ── LLMs para código ────────────────────────────────────────
\bibitem{codellama}
  \todo{CodeLlama / GPT-4 code review paper.}
\bibitem{llm_bugs}
  \todo{Trabalho sobre LLM para detecção de bugs.}
\bibitem{llm_hallucination}
  \todo{Trabalho sobre alucinações de LLM em code review.}

% ── Híbridos LLM + formal ───────────────────────────────────
\bibitem{llm_formal_hybrid}
  \todo{Trabalho sobre combinação LLM + verificação formal.}

% ── Datasets de bugs Python ─────────────────────────────────
\bibitem{bugsinpy}
  \todo{BugsInPy ou dataset similar de bugs Python reais.}

% ── Análise estática Python ─────────────────────────────────
\bibitem{pysa}
  \todo{Pyre, mypy, Pyflakes ou similar.}

\end{thebibliography}

\end{document}
